\section{大规模训练：数值故障、混淆因素与 R\&D 冰山}

当训练进入数千 GPU 和数周运行的规模，小实验的成功不再足以证明配置安全。课堂举例说，某个 FP16 路径在较小模型上看起来与 FP32 一致，却在大模型训练约两周后出现 loss 上升。这里最难的不是“知道 FP16 可能不稳定”，而是故障发生时同时存在数据、优化器、网络、硬件、代码和随机性等大量 \term{confounding factors}（混淆因素）。

\subsection{为什么小规模验证会漏掉大规模错误}

\term{FP16} 是 16-bit floating point；它节省显存和带宽，但指数和尾数范围有限。误差可能在单层、短序列或短训练中很小，却在更深网络、更长 horizon 和更高梯度方差下逐步放大。安全做法不是全面退回 FP32，而是识别敏感算子，使用 BF16、FP32 accumulation、loss scaling、数值监控和分阶段规模验证。

\lecturefigure{06-scale-amplifies-bugs.png}{小规模通过的数值近似，可能在长时间大训练中变成延迟故障。}{本地字幕 12:08--13:35；概念重绘。}

\begin{knowledgebox}{读图：延迟暴露会改变实验设计}
如果错误两周后才显现，单次二分排查就可能消耗数千 GPU-day。团队需要在训练早期记录梯度范数、激活范围、loss 分桶、硬件错误、数据 shard 和代码版本，并设计能尽早放大风险的 medium-scale canary。可观测性不是训练完成后的报表，而是降低调试搜索空间的实验仪器。
\end{knowledgebox}

\subsection{公开 final run cost 为什么会误导}

论文通常报告最终训练所用 GPU 数、tokens 和天数；这些数字只描述成功的主跑，不包含配方试验、失败 checkpoint、数据清洗、内核调优和系统搭建。课堂用 LLaMA 与 DeepSeek 的成本讨论提醒：把公开算力乘以云价，只能得到某次运行的近似账单，不能反推出团队获得该配方的总 R\&D 成本。

\lecturefigure{07-rd-iceberg.png}{最终成功训练只是可见成本，真正昂贵的是获得稳定配方的全过程。}{本地字幕 14:20--17:05；讲者成本估计仅作课堂背景。}

\begin{warningbox}{不要把讲者估计当成审计数据}
访谈中的 GPU 数量、小时单价和六个月总成本是现场口述，字幕还可能混淆 A100 与 H100。讲义保留“规模很大、试错很贵、final run 低估 R\&D”这一可教学结论，不用这些估计计算精确财务结论。
\end{warningbox}

\subsection{组织记忆：为什么“做过一次”很重要}

大训练 know-how 包含许多难以从论文复原的操作知识：什么指标先报警，哪类 loss spike 可以恢复，如何切分数据，何时重启，哪些优化只在特定拓扑有效。Lample 的创业判断是，团队已经在大厂支付过数月 trial-and-error 成本，因此能在资源更紧的创业环境中更快得到可靠训练。

\lecturefigure{08-data-heavy-team.png}{Mistral 早期团队把大部分人力放在数据，而不是只写训练代码。}{本地字幕 17:10--18:10；“1 人训练、6 人数据”为课堂口述的组织比例。}

\begin{knowledgebox}{读图：data work 为什么容易被低估}
数据工作包括来源审计、去重、质量过滤、语言配比、污染检测、困难样本分析和 evaluation slice 设计。它没有一个醒目的新 architecture 名称，却直接决定训练信号。老师甚至用“alienating”描述这种工作不够光鲜，但随后强调它最关键；这正是必须保留的 teacher voice。
\end{knowledgebox}

\teachervoice{课堂还比较了不同大厂基础设施的可迁移性：更高层的内部平台让日常工作方便，却可能隐藏调度、存储和训练栈细节；较低层环境要求研究者亲手处理更多问题，离开组织后反而更容易重建系统。这不是鼓励重复造轮子，而是提醒团队保留关键层的可解释性和操作能力。}

\subsection{本章小结}

大规模训练的困难来自新数值区域、昂贵反馈周期和知识隐性化。可靠团队不仅有更多 GPU，还需要 canary、可观测性、版本记录、数据能力和已经支付过的试错经验。

